\section{Split the reduction without changing its meaning}
A tile is a bounded interval in each relevant index. For tile edge \(T\), choose
starts \(I,J,L\) for rows, columns and reduction positions. Their exclusive ends are
\[
 I_e=\min(I+T,M),\qquad J_e=\min(J+T,N),\qquad L_e=\min(L+T,K).
\]
Process \(I\leq i<I_e\), \(J\leq j<J_e\), and \(L\leq k<L_e\).
The last tile can be shorter than \(T\) in any of the three dimensions.
Compute and clip each end independently. A correct row bound does not protect a
right-input column or a reduction position.

The illustrated order is \(I,J,L,i,j,k\). For one output tile, process all
reduction tiles before moving to the next output tile. Initialize every output
cell to zero before these loops begin. At the start of a reduction tile, reload
the cell's existing partial sum, add its next bounded group of products, and store
that new partial sum. Resetting to zero at every reduction tile discards all earlier
groups and computes only the final tile's contribution.

For this schedule the output invariant before reduction tile \(L\) is
\[
 C_{i,j}=\sum_{q=0}^{L-1}A_{i,q}B_{q,j}
 \quad\text{for the current output tile.}
\]
At \(L=0\), explicit initialization establishes zero. At the next tile start,
\(L\) is advanced by \(T\). The last clipped tile ends at \(K\), so its stored
value is the full dot product. All additions are exact integers in this exercise;
regrouping them introduces no floating-point rounding change.

\begin{MARSWorkedExample}{Keep the first reduction tile}
Use the two-by-three example with \(T=2\). After reduction positions zero and one,
its one output tile contains
\[
 \begin{pmatrix}25&28\\73&82\end{pmatrix}.
\]
The last reduction tile contains only position two. Its contributions are
\(33,36,66,72\), giving the final \(58,64,139,154\). A program that clears the
accumulator for the last tile stores only those four contributions. A program that
assumes two products in the last tile reads beyond the active input.
\end{MARSWorkedExample}

\begin{MARSExercise}{Clip all three ends}
For \((M,K,N,T)=(5,7,3,3)\), list the row, column and reduction intervals. How many
output tiles exist, and how many reduction tiles contribute to each? Check the
last element's address in each interval against the active byte length.
\end{MARSExercise}
